\documentclass[11pt]{article}
\usepackage[T1]{fontenc}
\usepackage[utf8]{inputenc}
\usepackage{lmodern}
\usepackage[margin=1in]{geometry}
\usepackage{amsmath,amssymb,amsthm,mathtools}
\usepackage{booktabs,array,longtable}
\usepackage{microtype}
\usepackage{xcolor}
\usepackage{listings}
\usepackage{enumitem}
\usepackage{hyperref}
\usepackage{bookmark}
\usepackage{fancyhdr}
\hypersetup{colorlinks=true,linkcolor=blue!45!black,urlcolor=blue!50!black,
 citecolor=blue!45!black,pdftitle={Exact Unknot Recognition: Budgeted Bracket Filters and a Compiled F2 Kernel},
 pdfauthor={OpenAI, prepared for Vladimir Reshetnikov}}
\pagestyle{fancy}\fancyhf{}
\fancyhead[L]{\small Exact unknot recognition: acceleration report}
\fancyhead[R]{\small 18 September 2026}
\fancyfoot[C]{\thepage}
\setlength{\headheight}{14pt}
\setlength{\parskip}{4pt}
\setlength{\parindent}{1.3em}
\setcounter{tocdepth}{2}
\lstset{basicstyle=\ttfamily\footnotesize,columns=fullflexible,breaklines=true,
 frame=single,framerule=.3pt,rulecolor=\color{black!25},showstringspaces=false,
 keepspaces=true,aboveskip=8pt,belowskip=8pt}
\newtheorem{theorem}{Theorem}[section]
\newtheorem{proposition}[theorem]{Proposition}
\newtheorem{lemma}[theorem]{Lemma}
\newtheorem{corollary}[theorem]{Corollary}
\theoremstyle{remark}\newtheorem{remark}[theorem]{Remark}
\newcommand{\F}{\mathbb F}
\newcommand{\Z}{\mathbb Z}
\newcommand{\Q}{\mathbb Q}
\newcommand{\Kh}{\operatorname{Kh}}
\newcommand{\rKh}{\widetilde{\operatorname{Kh}}}
\newcommand{\End}{\operatorname{End}}
\newcommand{\code}[1]{\texttt{\detokenize{#1}}}
\title{\textbf{Exact Unknot Recognition with Budgeted Bracket Filters\\
 and a Compiled $\F_2$ Cobordism Kernel}\\[8pt]
 \large An implemented and tested acceleration of \code{Knots.zip/fast/}}
\author{Prepared for Vladimir Reshetnikov\\\small Computational research and implementation: OpenAI}
\date{18 September 2026 \quad---\quad Software version 0.2.0}
\begin{document}
\maketitle
\begin{abstract}
We supply a working replacement for the archive's exact unknot recognizer, not a
quasi-polynomial claim or a hierarchy-algorithm stub. The main practical change
is to avoid constructing Khovanov homology when an inexpensive, exact invariant
obstruction already proves knottedness. A sparse modular determinant and a
budgeted Kauffman-bracket scan are followed, when appropriate, by an exact
integer bracket evaluation; inconclusive tests always fall through to the
complete backend. Inside that backend, topology compilation, characteristic-two
algebra, bounded memoization, and an equivalent heap-based scan planner reduce
work without changing the computed ranks. The planner improves from
$O(rn^2)$ to $O(rn\log n)$ for $r$ starting points, and the descending test becomes
linear. We also correct the writhe convention for general valid PD inputs.
In 138 paired cold-process trials, the median full-recognition speedup is
$568.54$ on a 22-crossing connected sum of two Conway knots; raw Khovanov
computation is $42.85$ times faster on a 1024-crossing unknot chain. Regressions
and unresolved stress cases are reported. Forty-two tests pass, including
independent state sums and reduced-cube calculations. The general worst-case
bound remains exponential. We identify precisely why small scan width does not
by itself make the explicit Khovanov backend quasi-polynomial, and separate the
implemented results from the additional obligations in Lackenby's hierarchy
approach.
\end{abstract}
\newpage
\begingroup
\small\setlength{\parskip}{0pt}
\tableofcontents
\endgroup
\newpage

\section{Scope, outcome, and reproducibility}
\label{sec:scope}
The target is the Python package in \code{fast/} of the supplied
\code{Knots.zip}~\cite{Input}. The original seven module files are preserved byte-for-byte
in \code{reference/baseline_fastunknot/}; they are imported under a different
package name solely to allow paired experiments. The modified package is
\code{fastunknot/}. It has no runtime dependency outside Python's standard
library and is designed for Python 3.10 or later. This delivery was executed
with Python 3.13.5.

The outcome has three distinct levels. First, the complete recognizer is
substantially faster on several difficult-for-Alexander inputs. Second, there
are proved improvements to individual algorithms and a proved width-sensitive
bound for the new one-sided bracket filter. Third, \emph{there is no general
$n^{O(\log n)}$ bound for the delivered recognizer}. The last point is not an
assumption about what future research can achieve: it describes the actual
implementation and the bounds established here.

The most important design constraint is that \code{UNKNOWN}, equality of an
invariant with the unknot's value, and exhaustion of a filter budget must never
be confused with a topological verdict. No floating-point computation is used
for a mathematical decision. The algorithm does not consult a knot table or
recognize benchmark fixtures by name.

\subsection{The retained input contract}
A crossing is a counterclockwise tuple $(a,b,c,d)$, with under-strand $a$--$c$
and over-strand $b$--$d$. Each edge label occurs twice. The validator normalizes
labels, checks that following opposite ports gives a single component, and
checks that the rotation system is spherical. For a connected four-valent
$n$-crossing map, the latter check is the Euler condition $F=n+2$.
The input \code{{"pd": []}} represents one crossing-free circle, not an
empty link. Links and virtual diagrams are rejected. The inherited braid and
grid input formats are retained. Complexity bounds below use normalized edge
labels; reading a noncanonical encoding additionally costs its input bit length. Grid rows run bottom to top and vertical
segments pass over horizontal segments.

\begin{lstlisting}
{"pd": [[1,4,2,5], [3,6,4,1], [5,2,6,3]]}
{"braid": {"strands": 3, "word": [1,-2,1,-2]}}
{"rows": [[0,2], [1,3], [2,4], [0,3], [1,4]]}
\end{lstlisting}

\subsection{What is new, and what is inherited}
The inherited recognizer validates, performs crossing-decreasing Reidemeister
I/II moves, tests for a descending presentation, computes the full Alexander
polynomial, and finally computes reduced Khovanov rank through an unreduced
Bar--Natan scan. We retain the underlying topology and elimination strategy.
The new work consists of executable filters, algebraic specialization and
caching, faster exact-equivalent planning, a linear descending test, convention
repairs, additional validation, a result verifier, and reproducible experiments.
The Kauffman bracket, Bar--Natan's local method, and Khovanov unknot detection
are established mathematics, not discoveries claimed by this report
\cite{Kauffman,BN,KM}.

\section{Lackenby's material and the quasi-polynomial target}
\label{sec:lackenby}
The attached talk is dated February 2021, whereas the attached preprint is
\emph{Incompressible surfaces, hierarchies and unknot recognition},
arXiv:2607.23350v1, dated 25 July 2026. These are different documents and should
not be conflated. The talk announces the quasi-polynomial result and describes
an accelerated construction involving hierarchical multi-surfaces, Cheeger
regions, cutting, weak reductions, and simplification. The preprint gives a
hierarchy-based algorithm and further structural results. In Section 9 it
explicitly separates an iteration count from a running-time estimate: its
initial count uses a maximal hierarchy length $L$ and maximal pattern
complexity $g$, with bound $L(g+1)^L$ \cite{LackenbyTalk,Lackenby2026}.

For clarity, the following is an elementary conditional calculation, not a
new proof of the announced topological algorithm.
\begin{proposition}[A sufficient complexity contract]
Suppose a constructive hierarchy algorithm on an $n$-crossing input has at most
$L(g+1)^L$ iterations, where $L\le C\log(n+2)$ and $g\le(n+2)^C$.
Suppose all represented data and every iteration can be computed using at most
$n^{O(\log n)}$ bit operations and comparable space. Then its running time is
$n^{O(\log n)}$.
\end{proposition}
\begin{proof}
The logarithm of the iteration bound is
$\log L+L\log(g+1)=O((\log(n+2))^2)$. Multiplying it by another
$\exp(O((\log(n+2))^2))$ bound leaves the same asymptotic class.
\end{proof}

A correct implementation of that route must supply the represented handle
structures and boundary patterns, exact surface operations, a decreasing
potential, size bounds, and the accelerated depth bound together. A finite
lexicographically decreasing counter is not sufficient: its length, digit
sizes, coordinate bit lengths, and the work required to find each next move
matter. Section 10 and the later sections of the preprint develop refinements;
our implementation does not purport to instantiate them or to turn the talk's
flowchart into executable primitives.

The delivered package instead makes genuine progress on the requested fallback
goal: better practical performance and better asymptotic subroutines for the
provided implementation. There are no unimplemented hierarchy functions hidden
behind the command-line interface. Section~\ref{sec:multiplicity} also exhibits
a concrete obstruction to a tempting but invalid alternative proof based only
on the number of frontier matchings.

\section{The accelerated recognition pipeline}
\label{sec:pipeline}
Write $n$ for the number of crossings after optional Reidemeister reduction
when discussing filter budgets. The actual pipeline is:
\begin{enumerate}[itemsep=1pt,topsep=3pt]
\item Validate the original diagram; optionally simplify by legal R1/R2 moves
and run the descending test. These stages may prove \code{UNKNOT}.
\item Evaluate an Alexander minor at $t=-1$ modulo $2^{31}-1$ using sparse
Gaussian elimination. A value other than $\pm1$ proves \code{KNOTTED}.
\item Compute a budgeted bracket value at $A=2$ modulo $2^{31}-1$. A mismatch
with the unknot value proves \code{KNOTTED}.
\item If that modular evaluation completed but matched, try an exact integer
scaled bracket evaluation at $A=2$, again with a budget. A mismatch proves
\code{KNOTTED}.
\item Compute the original full exact Alexander polynomial, unless disabled.
A nonunit normalized polynomial proves \code{KNOTTED}.
\item Run the improved complete Khovanov backend. Reduced rank one proves
\code{UNKNOT}; any other rank proves \code{KNOTTED}.
\end{enumerate}

Each bracket evaluation receives a default transition budget
\begin{equation}
 B(n)=10000+200n
 \label{eq:budget}
\end{equation}
and a limit of 4096 live pairing states. A transition means one smoothing of
one current pairing. Reaching either ceiling discards the incomplete test and
continues the pipeline. It does not consume the right to try the complete
backend. The second, exact evaluation is skipped when the first did not
complete, to avoid predictably repeating an already oversized computation.

The determinant precedes the bracket because it often rejects nontrivial
knots at very low cost. The bracket precedes the full polynomial because it
handles the Alexander-trivial fixtures and their connected sums particularly
well. This order is a practical choice, not a theorem of optimality. It can be
slower on small inputs whose Alexander polynomial already rejects them; the
measurements below include these cases.

\begin{theorem}[Soundness and completeness of the wrapper]
On a validated classical one-component input, assuming the stated exact
algebraic operations and the retained topological algorithms, every returned
\code{UNKNOT} or \code{KNOTTED} is correct. Without an overall time or
Khovanov-object ceiling, the recognizer is complete. The default finite
budgets for the new filters do not affect completeness.
\end{theorem}
\begin{proof}
R1/R2 moves preserve knot type. A descending closed diagram represents the
unknot: cut at its descending basepoint, place the successive traversed arc at
strictly decreasing height, and then straighten the resulting unentangled arc
before closing it. Sections~\ref{sec:det} and \ref{sec:bracket} prove that each
new rejection is a necessary-identity violation for the unknot. A completed
full Alexander calculation is an established invariant obstruction.
Every inconclusive filter leaves the same reduced diagram for the next stage.
The retained Khovanov construction is finite, and the rank-one detection
criterion is justified in Section~\ref{sec:detection}. With imposed resource
ceilings an unfinished complete calculation returns \code{UNKNOWN}, not a
verdict.
\end{proof}

\section{A sparse modular determinant obstruction}
\label{sec:det}
Collapse the over-strand edges at each crossing to obtain the Wirtinger arcs.
For a nonempty one-component $n$-crossing diagram there are $n$ such arcs. An
Alexander matrix has one row per crossing; up to the crossing convention its
three contributions are $1-t$ on the over-arc and $t,-1$ on the two under-arcs.
At $t=-1$, the two under entries agree, so crossing signs and their incoming
order disappear. The row becomes
\begin{equation}
 2e_{\mathrm{over}}-e_{\mathrm{under,1}}-e_{\mathrm{under,2}}.
 \label{eq:foxminusone}
\end{equation}
Contributions are added when two arc indices coincide. Delete one row and one
column. The resulting minor is $\Delta_K(-1)$ up to sign, because the usual
unit ambiguity $\pm t^j$ becomes $\pm1$.

\begin{proposition}
Let $p=2147483647$. If the determinant of this minor modulo $p$ lies outside
$\{1,p-1\}$, the knot is nontrivial.
\end{proposition}
\begin{proof}
The Alexander polynomial of the unknot is one up to the unit ambiguity.
Therefore every indicated minor for an unknot specializes to $+1$ or $-1$.
Reduction modulo $p$ preserves this necessary identity. Its violation is an
exact disproof, not a probabilistic guess.
\end{proof}

Rows are sparse Python dictionaries. Gaussian elimination chooses the first
available nonzero entry in each fixed pivot column, swaps rows when necessary,
multiplies the running determinant by the pivot, and updates only stored tail
entries. Modular inverses are computed exactly. In the dense worst case this
uses $O(n^3)$ field operations. The recognizer uses only the fixed prime above;
the low-level determinant routine assumes its optional modulus is prime and
is not a general composite-modulus determinant routine.

A determinant congruent to $\pm1$ is inconclusive. This includes every
Alexander-trivial knot. The full Alexander polynomial is retained because
$\Delta_K\ne1$ can hold even when the determinant test does not reject.

\section{Bracket filtering by frontier aggregation}
\label{sec:bracket}
\subsection{Conventions and the necessary identity}
For $A$ invertible, put
\begin{equation}
 \delta=-A^2-A^{-2}.
\end{equation}
We use the unnormalized bracket $\mathcal B_D(A)$: every closed state circle
contributes $\delta$, including the first. At a PD crossing $(a,b,c,d)$, the
zero smoothing pairs $(a,b),(c,d)$ with weight $A$, and the one smoothing pairs
$(a,d),(b,c)$ with weight $A^{-1}$. Thus, for $n>0$,
\begin{equation}
 \mathcal B_D(A)=\sum_{s\in\{0,1\}^n}A^{n-2|s|}\delta^{c(s)},
 \label{eq:statesum}
\end{equation}
where $c(s)$ is the number of circles of the completely smoothed diagram.
The zero-crossing input has bracket $\delta$.

With the oriented writhe $w(D)$ in the PD convention implemented here, an
unknot diagram must satisfy
\begin{equation}
 \mathcal B_D(A)=\delta(-A^3)^{w(D)}.
 \label{eq:unknotbracket}
\end{equation}
The power is computed in the ring in question; a negative writhe is handled by
inverting the unit $-A^3$. In particular, the modular test does not divide by
$\delta$, which might be zero.

For completeness, the local algebra underlying this identity can be checked
without a knot table. Let $e_i$ be the Temperley--Lieb smoothing, satisfying
$e_i^2=\delta e_i$ and $e_ie_{i+1}e_i=e_i$. A crossing and its inverse have
operators $AI+A^{-1}e_i$ and $A^{-1}I+Ae_i$; their product is
$I+(A^2+A^{-2}+\delta)e_i=I$, proving R2 invariance. Expanding the adjacent
three-crossing products, their difference is a scalar multiple of
$A^2+\delta+A^{-2}=0$, proving R3 invariance. The two R1 curls evaluate to
$-A^3$ and $-A^{-3}$. Multiplication by $(-A^3)^{-w(D)}$ removes these factors.
The crossing-free circle then proves~\eqref{eq:unknotbracket}. This is the
standard bracket construction \cite{Kauffman}.

\subsection{The state and transition}
After some crossings have been processed, keep the edges incident to exactly
one processed crossing as the frontier. A partial smoothing consists of
closed circles and paths joining frontier endpoints. Closed circles are
immediately evaluated. The remaining information is a perfect matching $m$
of the frontier and a coefficient $v(m)$ in the chosen ring.

To add a crossing, take the union of the current matching and one smoothing's
two local arcs. This is a degree-at-most-two multigraph: its path components
give a new matching $m'$, and its cycle components give a number $c$ of newly
closed circles. Add
\begin{equation}
 v(m)A^{1-2s}\delta^c
 \quad\hbox{to the coefficient of }m',\qquad s\in\{0,1\}.
 \label{eq:transition}
\end{equation}
An edge may occur twice in a PD row. Degrees and edges are counted with
multiplicity, so a self-loop smoothing is correctly recognized as a circle.
The function \code{attach} implements this transition independently of the
Khovanov backend's gluing routine.

\begin{lemma}[Exact aggregation]
After every processed prefix, the coefficient stored at $m$ is the sum of the
weights of all partial smoothing choices with residual matching $m$, with
closed circles already evaluated. Consequently the completed scan equals
\eqref{eq:statesum} for every crossing order.
\end{lemma}
\begin{proof}
Initially there is one empty partial choice with weight one. Given the
invariant for a prefix, each choice extends uniquely in two ways at the next
crossing. Equation~\eqref{eq:transition} contributes exactly the new crossing
weight and the newly completed circles. Future attachments depend only on the
residual pairing, not on how it arose, so equal pairings may be summed.
At the end all components are closed and the only possible matching is empty.
The zero-crossing circle is treated separately.
\end{proof}

In modular mode, coefficients are reduced modulo $p$ after every addition.
Terms whose aggregate coefficient vanishes are omitted. This omission is safe
by linearity. A completed mismatch in~\eqref{eq:unknotbracket} returns a
knottedness witness. A matching value proves nothing about triviality.

\subsection{Exact integer mode without rational denominators}
A fixed modular value can collide with the unknot value. To avoid relying only
on finite-field information, the implementation also has an exact evaluation
at any integer $a\ge2$, normally $a=2$. Put $D=-(a^4+1)$, so
$\delta=D/a^2$. Instead of storing rational coefficients after $j$ crossings,
store them multiplied by $a^{5j}$. The transition weight becomes
\begin{equation}
 a^5a^{1-2s}(D/a^2)^c=a^{6-2s-2c}D^c.
 \label{eq:integertransition}
\end{equation}
At most two circles close when the two smoothing arcs are attached; hence
$c\le2$, and the exponent in~\eqref{eq:integertransition} is nonnegative.
All intermediate arithmetic is integral.

For $n\ge1$ the expected scaled unknot value is
\begin{equation}
 U_D=(-1)^{w(D)}D\,a^{5n+3w(D)-2}.
 \label{eq:integertarget}
\end{equation}
Since $|w(D)|\le n$, its exponent is at least $2n-2$. For $n=0$ the scale is
$a^2$ and both values equal $D$. Completed values are serialized as hexadecimal
strings, avoiding decimal digit-conversion limits on very large Python
integers. These are exact integer comparisons, not numerical approximations.

For fixed $a$, each transition multiplier has bounded absolute size, and at
most $2^n$ terms contribute. Therefore every coefficient has $O(n)$ bits.
For variable $a$, a bound is $O(n(1+\log a))$ bits. An equality at one exact
point is still only inconclusive: the algorithm does not assume that the
Jones polynomial, or any one of its evaluations, detects the unknot.

\subsection{Width-sensitive and budget-sensitive complexity}
Let $w$ be the largest frontier size of the supplied order, and define
\begin{equation}
 M(w)=(w-1)!!=\frac{w!}{2^{w/2}(w/2)!}\quad(w\text{ even}),\qquad M(0)=1.
 \label{eq:matchingbound}
\end{equation}
There are at most $M(w)$ perfect matchings at a step. We use this safe bound,
not a Catalan bound that would require an additional disk-boundary embedding
hypothesis. An arbitrary processed region need not be a disk.

\begin{proposition}[Complexity of the bracket filter]
For fixed modulus and evaluation point, an uncapped scan of width $w$ uses
$O(nM(w)\operatorname{poly}(w))$ word operations, apart from order construction.
It stores $O(M(w)\operatorname{poly}(w))$ words. Since
$M(w)=\exp(O(w\log(w+1)))$, this is a width-sensitive bound for the filter.
At a fixed integer point, a conservative bit bound is
$\operatorname{poly}(n,w)M(w)$. With the default transition budget and state
ceiling, the added filter work is polynomial in $n$ even for unbounded width.
\end{proposition}
\begin{proof}
Each state has two transitions per crossing. The degree-two connectivity
calculation, canonical sorting, and hashing have polynomial cost in $w$;
$O(w^2)$ is a sufficient conservative bound for the simple implementation on
word-sized labels. Modular scalar operations have fixed-size operands.
The integer case additionally uses the linear coefficient bit bound above.
Under the default budget there are only $O(n)$ transitions in each of at most
two evaluations, and $w\le2n$. The matching keys and the exact coefficients
therefore have polynomial size. Crossing-order construction is polynomial as
well. Interrupting the scan cannot increase its work beyond these caps except
for checking and constructing the bounded current transition.
\end{proof}

This theorem is deliberately not stated for the Khovanov complex: its object
multiplicities are additional data, as Section~\ref{sec:multiplicity} shows.

\section{A necessary orientation repair}
\label{sec:writhe}
The old \code{Diagram.signs()} reads which over-port is incoming but implicitly
assumes that the incoming under-port is slot zero. That assumption is not part
of the accepted PD format. Rotating a row by two positions preserves the
crossing, but can invalidate that implicit orientation choice.

For each crossing let $u\in\{0,2\}$ and $o\in\{1,3\}$ be the incoming under-
and over-port indices along an oriented traversal. The corrected sign is
\begin{equation}
 \epsilon=\begin{cases}+1,&o-u\equiv3\pmod4,\\-1,&o-u\equiv1\pmod4.
 \end{cases}
\end{equation}
Reversing the knot orientation adds two to both incoming indices, so the sign
is unchanged. Rotating a PD row by two has the same invariance. A mirror
reverses the sign. This is essential for~\eqref{eq:unknotbracket}; a writhe error
would turn a correct bracket computation into a false rejection.

For example, the three-crossing braid \code{[1,1,1]} is converted by this
package to
\[
 (5,4,0,1),\quad(1,0,2,3),\quad(3,2,4,5).
\]
Rotating the first row by two leaves the new writhe equal to $-3$, whereas the
old function changes from $-3$ to $+1$. The signs here reflect the package's PD
and braid conversion convention; the example is not an assertion about a
different external braid-sign convention.

The original Alexander code always assigns its under-arc variables using PD
slots zero and two. Its old sign convention is compensated by that assignment.
Accordingly, this report does \emph{not} infer that the old Alexander polynomial
or the old recognizer was wrong from the writhe bug alone. The new Alexander
matrix uses the genuine oriented sign \emph{and} swaps the incoming/outgoing
under variables when the incoming under-port is slot two. Exact comparisons
on rotated inputs confirm preservation of the polynomial.

\section{Faster exact-equivalent preprocessing and planning}
\subsection{A linear descending test}
Fix one orientation and index its $2n$ crossing visits cyclically. If a crossing
is visited over at index $o$ and under at index $u$, precisely the start indices
in the circular interval $(o,u]$ encounter its underpass first. Add one on each
such bad interval with a difference array, splitting a wrapped interval into
two ordinary intervals. Prefix sums give the number of bad crossings for every
basepoint. A zero is a descending presentation. Repeat for the reverse
orientation and choose the smallest valid dart, matching the old tie-breaking
rule.

There are $n$ intervals and $2n$ starts in each of two orientations. The running
time and auxiliary space are $O(n)$, compared with up to $4n$ separate
$O(n)$ traversals in the old test. This optimization does not alter the R1/R2
simplifier; that simplifier still recomputes faces and validates diagrams after
moves, and remains a polynomial but potentially significant cost.

\subsection{Heap-based scan ordering}
For an unprocessed crossing $i$, let $s_i$ count its edges shared with the
current processed region, let $\ell_i$ count self-loop pairs at that crossing,
and let $b$ be the current boundary size. The old greedy score is
\[
 \bigl(s_i,-(b+4-2s_i-2\ell_i),-i\bigr).
\]
For a fixed step its ordering is exactly the ordering of
$(s_i,\ell_i,-i)$. Self-loop counts are fixed; only shared counts change.

Build edge-to-crossing incidence lists. Store crossings in a heap keyed by
$(-s_i,-\ell_i,i)$, using the opposite sign for a minimum heap. When a crossing
is processed, increment the shared count of each still-active neighboring
crossing and push its new key. Discard inactive or stale heap entries on pop.
Each diagram edge causes at most one such increment. Thus there are $O(n)$
heap operations, each costing $O(\log n)$.

\begin{proposition}
For the same starting crossing, the new planner returns exactly the original
order in $O(n\log n)$ time instead of $O(n^2)$. Trying $r$ starts gives
$O(rn\log n)$ instead of $O(rn^2)$.
\end{proposition}
\begin{proof}
The score reduction preserves the total order, including index tie-breaks.
The heap contains a current key for every active crossing; stale keys cannot
cause an incorrect choice because they are rejected. Incidence updates give
exactly the number of processed neighbors, with multiplicities. Induction on
the chosen prefix proves equality of the orders. The operation count gives
the complexity bound.
\end{proof}

The default start-selection rule and width-profile comparison are unchanged.
The old stride convention may produce slightly more than the requested number
of starts, but still $O(r)$. Requesting all starts is $O(n^2\log n)$ in the new
version, not $O(n\log n)$. The benchmarks use the bounded-start policy of the
recognizer. This is not an optimal-width algorithm or a separator algorithm.

\section{Compiling the characteristic-two cobordism algebra}
\label{sec:kernel}
\subsection{Representation and cancellation}
The scan is Bar--Natan's local simplification method \cite{BN}. Objects are
frontier matchings, with homological degree and multiplicity. Closed circles
are delooped into two summands corresponding to $1,x$ in
$\F_2[x]/(x^2)$. A morphism is an $\F_2$ sum of dot-subsets on the boundary
circles of a cobordism; addition is symmetric difference of these finite sets.
Invertible differential entries between equal matchings are cancelled after
each crossing.

If $\phi:B\to C$ is an invertible differential block, with neighboring blocks
$\delta:A\to C$ and $\gamma:B\to F$, Gaussian cancellation replaces the
$A\to F$ block by
\begin{equation}
 d_{AF}+\gamma\phi^{-1}\delta.
 \label{eq:cancel}
\end{equation}
The usual minus sign equals plus in characteristic two. The eliminated pair is
contractible, so this preserves homotopy type. We retain the old LIFO pivot
policy; improvements in the measurements are not attributed to an unimplemented
minimum-fill pivot heuristic.

\subsection{Every unit is its own inverse}
If a matching has $r$ arcs, its endomorphism ring has the form
\begin{equation}
 R_r=\F_2[x_1,\ldots,x_r]/(x_1^2,\ldots,x_r^2).
\end{equation}
Its units are exactly $1+\nu$ with zero constant term in $\nu$.
\begin{lemma}
Every such unit satisfies $(1+\nu)^{-1}=1+\nu$.
\end{lemma}
\begin{proof}
In a commutative characteristic-two ring the square of a sum is the sum of its
squares. Every nonconstant square-free monomial has zero square, because it
contains some $x_i$ and $x_i^2=0$. Hence $\nu^2=0$ and
$(1+\nu)^2=1$.
\end{proof}

The new inverse therefore validates the constant term and copies the set of
terms. Its cost is linear in the stored representation, rather than requiring
morphism multiplications. Importantly, the old geometric-series loop already
terminates after at most two products in this ring: we do not claim that it
required $r$ or exponentially many iterations. The improvement is eliminating
these expensive products altogether.

For endomorphism multiplication, square-free monomials indexed by dot sets
$S,T$ satisfy
\[
 x_Sx_T=0\quad(S\cap T\ne\varnothing),\qquad
 x_Sx_T=x_{S\cup T}\quad(S\cap T=\varnothing).
\]
The optimized path uses this rule directly and XORs equal output monomials.
Identity multiplication is also short-circuited.

\subsection{Separating topology from dot data}
For a general composite $a\to b\to c$, the topology depends on the three
matchings, not on the coefficient terms. The new \code{_composition_plan}
compiles once the connected components of the glued surface, their Euler
characteristics, their outgoing boundary-circle sets, and maps from each
incoming dot position to a component bit.

For each pair of input terms, the implementation only has to form the mask of
dotted components. A repeated bit means two dots lie on the same component and
the contribution is zero. A component of genus at least one is also zero:
its handle map is multiplication by $2x=0$ over $\F_2$. For a planar component
with $q>0$ outgoing circles, zero dots gives the sum of the $q$ monomials in
which exactly one outgoing circle is undotted; one dot gives the all-dotted
monomial. A closed planar component evaluates to one precisely when it has
one dot. These are exactly the multiplication, comultiplication, and counit
rules of the same Frobenius algebra. The optimized expansion changes their
organization, not their values.

New morphism and crossing-entry caches store frozen values and return copies
where the public internal API expects mutable sets. This prevents a caller's
subsequent XOR or clearing operation from poisoning a cached morphism. All
new caches have explicit LRU entry ceilings. The existing circles/glue caches
also receive finite entry ceilings, and \code{clear_caches()} resets them.
These limits bound the number of entries, \emph{not} the byte size of a single
large entry. No fixed memory bound in megabytes is claimed.

The limits are 32768 for circles and glue, 8192 for composition plans, 32768
for mask expansions, 16384 for completed compositions, and 8192 for crossing
entries. They are engineering choices, not complexity constants required by
correctness. Eviction merely causes recomputation.

\subsection{Why the final rank detects the unknot}
\label{sec:detection}
The backend computes the total unreduced rank over $\F_2$ and divides by two.
Here is the algebraic justification for that division. On each resolved state
let $X$ multiply the basepoint-circle factor by $x$, and let $\nu$ be the sum
of the operators $\partial_x$ on all circle factors, where
$\partial_x(1)=0$ and $\partial_x(x)=1$. Over $\F_2$ both are chain maps:
this follows directly by checking multiplication and comultiplication on the
basis $1,x$. They satisfy
\[
 X^2=\nu^2=0,\qquad X\nu+\nu X=1.
\]
The reduced subcomplex is $\operatorname{im}X$. Every chain splits uniquely as
$X\nu c+\nu Xc$, giving
\[
 C\cong\operatorname{im}X\oplus\nu(\operatorname{im}X),
\]
with the two summands isomorphic by $X$ and $\nu$ after suppressing the quantum
shift. Thus total unreduced rank is twice total reduced rank; see also the
characteristic-two structure discussed in \cite{Shumakovitch}.

Kronheimer--Mrowka prove that reduced rational Khovanov rank one detects the
unknot \cite{KM}. The integral reduced complex is free at chain level, so the
universal coefficient theorem gives
\[
 \dim_{\F_2}\rKh(K;\F_2)\ge\dim_{\Q}\rKh(K;\Q).
\]
The rational rank is at least one: the reduced graded Euler characteristic
specialized at one is one for a knot. Therefore reduced $\F_2$ rank one forces
rational rank one and hence triviality. Conversely, the unknot has reduced
rank one over $\F_2$. This proves the exact terminal criterion used here.

\section{Why frontier width alone cannot bound the explicit backend}
\label{sec:multiplicity}
The original README correctly warns in one place that matching multiplicities
need control, but elsewhere describes the exponential part as depending on
width rather than on crossing number. For the explicit minimal-complex
representation, a uniform $\operatorname{poly}(n)f(w)$ bound is false.

Let $C$ be the supplied Conway diagram and let $C_k$ be the connected sum of
$k$ copies, constructed by \code{tools/families.py}. It has $11k$ crossings.
The independent reduced cube and both scanning backends compute
\[
 \dim\rKh(C;\F_2)=33.
\]
For connected sums the reduced complexes tensor over the ground field, with
the marked-circle factors identified. On state bases this identifies the
unmarked factors of the two diagrams; crossing differentials act on their
own factor, giving the tensor differential. The K\"unneth theorem over a
field consequently gives
\begin{equation}
 \dim\rKh(C_k;\F_2)=33^k.
 \label{eq:rankpowers}
\end{equation}
This is the standard reduced connected-sum construction \cite{KhovanovPatterns}.
For $k=2$, the delivered tests and baseline benchmark independently of the
filter path verify reduced rank $1089$ and unreduced rank $2178$.

The factors can be scanned consecutively with bounded frontier size: between
fixed-size factors only the connecting strands are exposed, and within a
factor only a constant number of additional edges occurs. Nevertheless the
fully reduced closed complex has $2\cdot33^k$ explicit objects. At that final
stage there is one possible matching---the empty matching---but an exponential
multiplicity. Counting possible matchings cannot bound those multiplicities.
This is an output-size lower bound for this representation, not a lower bound
for every possible unknot recognizer or a compressed rank algorithm.

There is a corresponding positive result for the new filter. The exact
normalized bracket evaluation at $A=2$ of the supplied Conway diagram is
\begin{equation}
 r=\frac{\mathcal B_C(2)}{\delta(-2^3)^{w(C)}}
   =\frac{970142522911}{65536}>1.
 \label{eq:conwaypoint}
\end{equation}
The normalized bracket is multiplicative under connected sum: joining the
marked state circles reduces their combined circle count by one, cancelling
one factor of $\delta$ from the product, and writhe is additive. Hence the
corresponding value for $C_k$ is $r^k\ne1$.

\begin{corollary}[A genuine family-specific separation]
With a factor-consecutive bounded-width order and sufficient transition/state
budgets, the implemented exact integer bracket filter rejects every $C_k$ in
polynomial bit time. An algorithm that explicitly materializes the closed
minimal Khovanov complex for the same family requires exponentially many
objects.
\end{corollary}
\begin{proof}
Equation~\eqref{eq:conwaypoint} and multiplicativity give a mismatch for every
$k\ge1$. Bounded width gives a constant matching-state bound; coefficients
have $O(k)$ bits. There are $O(k)$ crossing steps, hence polynomial bit work.
Equation~\eqref{eq:rankpowers} proves the explicit-object lower bound.
\end{proof}

This corollary does not assert that the default greedy order and default caps
have the same bound for every $k$. It describes the exposed low-level filter
with the indicated order and budgets. The family tests check exact
multiplicativity for $k=1,\ldots,5$. A fixed modular evaluation alone would not
justify the corollary: in a finite field nonzero values have finite
multiplicative order, so powers can collide with one. That is a concrete reason
for including exact integer mode.

The global worst-case upper bound of the complete recognizer remains
$2^{O(n)}$, with conservative polynomial factors absorbed into the exponent.
The ordinary resolution cube has exponentially many states and enhanced-state
basis elements; intermediate delooping and finite-field elimination can be
bounded by exponential-size matrices and exponentially many dot terms.
The polynomially capped new filters do not alter that class. A general
quasi-polynomial guarantee would require an additional structural algorithm,
not just these optimizations.

\section{Paired measurements, including regressions}
\label{sec:bench}
\subsection{Protocol}
\code{benchmarks/run.py} generated 23 cases and ran each engine three times,
for 138 separate timed calls. The original ten fixtures are included, together
with connected-sum stress cases and raw-backend/planning diagnostics. Every
call uses a fresh subprocess and cold caches. Engine order alternates across
cases and repetitions. The child interpreter uses \code{-S} to avoid loading
unrelated site initialization; both packages need only the standard library.
\code{PYTHONHASHSEED=0} is fixed. The timer is \code{time.perf_counter} around
the requested function, excluding interpreter/import startup and the first
\code{Diagram.from_json} parse. Validation performed again inside the new
recognizer or raw backend is included in its time.

The environment reports Python 3.13.5, GCC 14.2.0, Linux x86-64 with glibc
2.41, and an AMD EPYC 9V74 80-Core Processor; five logical CPUs are visible to
the environment. The computation is serial. These are shared-container wall
clock measurements, not controlled hardware cycle counts or statistical
confidence intervals. Three trials provide a reproducible sample, not a claim
of universal speedup.

The cooperative budget is three seconds and the parent's hard subprocess
limit is 5.5 seconds. The original raw backend starts its internal deadline
\emph{after} order selection; this is why a completed original 1024-crossing
call can take more than three seconds in total. The new backend starts before
planning. These soft timers are not exact process deadlines. All 138 calls
returned; 126 completed, and 12 returned resource-limit results. There were no
hard timeouts or execution errors in this run.

The JSON contains every input, option, result, trial time, outcome, and
reported peak resident-set value. RSS units are platform-specific and it
includes the process environment, so it is not used as a memory-speedup claim.
The CSV includes minima, maxima, and medians. Table source is generated from
that JSON by \code{benchmarks/make_tables.py}, rather than transcribed by hand.

\begin{table}[htbp]\centering\small
\begin{tabular}{@{}lrrrr@{}}
\toprule Case & $n$ & Original (ms) & New (ms) & Ratio \\
\midrule
Crossing-free unknot & 0 & 0.020 & 0.026 & 0.75 \\
Trefoil & 3 & 0.092 & 0.123 & 0.75 \\
Figure-eight & 4 & 0.134 & 0.128 & 1.05 \\
Hard unknot (8 crossings) & 8 & 1.629 & 2.242 & 0.73 \\
Kinoshita--Terasaka & 11 & 39.313 & 1.018 & 38.61 \\
Conway & 11 & 39.304 & 0.822 & 47.79 \\
$T(3,5)$ & 10 & 0.591 & 0.912 & 0.65 \\
Unknot braid (40 crossings) & 40 & 6.681 & 7.999 & 0.84 \\
Determinant-one grid knot & 10 & 0.721 & 0.957 & 0.75 \\
Scrambled grid unknot & 6 & 0.259 & 0.288 & 0.90 \\
Conway $\#\,$ Conway & 22 & 1399.518 & 2.462 & 568.54 \\
Three Conway summands & 33 & \textsc{limit} & 3.379 & -- \\
Four Conway summands & 44 & \textsc{limit} & 4.112 & -- \\
\bottomrule\end{tabular}
\caption{Full recognition, defaults in both versions. Ratio is original/new; values below one are regressions. Each time is the median of three completed cold-process calls. Limit means all three calls returned UNKNOWN.}\label{tab:recognition}
\end{table}

\begin{table}[htbp]\centering\small
\begin{tabular}{@{}lrrrr@{}}
\toprule Case & $n$ & Original (ms) & New (ms) & Ratio \\
\midrule
Conway & 11 & 35.973 & 13.269 & 2.71 \\
Kinoshita--Terasaka & 11 & 39.206 & 15.090 & 2.60 \\
Hard unknot (8 crossings) & 8 & 2.476 & 2.165 & 1.14 \\
$T(3,5)$ & 10 & 10.206 & 9.958 & 1.02 \\
Unknot chain (256 crossings) & 256 & 231.987 & 20.751 & 11.18 \\
Unknot chain (1024 crossings) & 1024 & 3601.461 & 84.039 & 42.85 \\
Random 5-braid (36 crossings) & 36 & \textsc{limit} & \textsc{limit} & -- \\
\bottomrule\end{tabular}
\caption{Raw Khovanov backend, without recognition preprocessing or invariant filters. Both engines use the same deterministic scan-order policy.}\label{tab:raw}
\end{table}

\begin{table}[htbp]\centering\small
\begin{tabular}{@{}lrrrr@{}}
\toprule Case & $n$ & Original (ms) & New (ms) & Ratio \\
\midrule
Order only: 256-crossing chain & 256 & 217.852 & 5.844 & 37.28 \\
Order only: 1024-crossing chain & 1024 & 3550.885 & 27.065 & 131.20 \\
No R1/R2/descending: hard unknot & 8 & 2.907 & 3.203 & 0.91 \\
\bottomrule\end{tabular}
\caption{Planning and forced-fallback diagnostics. In the last row only Reidemeister reduction and descending recognition are disabled; invariant filters are still enabled.}\label{tab:aux}
\end{table}

\clearpage

\subsection{What the results establish}
The strongest completed end-to-end result is the two-Conway connected sum:
$1.399518$ seconds becomes $0.002462$ seconds, a median ratio of $568.54$.
This compares recognition with recognition, not an unfiltered old homology
calculation against a new shortcut with different requested output. Both
versions return \code{KNOTTED}. The baseline reaches reduced Khovanov rank
$1089$; the new version proves a bracket mismatch before constructing that
complex.

Conway and Kinoshita--Terasaka each have trivial Alexander polynomial in these
fixtures. The full recognizer is respectively $47.79$ and $38.61$ times faster.
The raw backend still speeds up by $2.71$ and $2.60$ on the same PD inputs,
showing that the implementation is not merely a wrapper around an unchanged
exponential kernel.

For the chain braid $\sigma_1\sigma_2\cdots\sigma_n$ on $n+1$ strands, raw
backend time changes from $0.231987$ to $0.020751$ seconds at $n=256$, and from
$3.601461$ to $0.084039$ seconds at $n=1024$. Order-only times isolate the
asymptotic planner change. These raw cases intentionally bypass Reidemeister
recognition: the chain is easy to simplify, so they are backend scaling tests,
not claims that an otherwise difficult unknot has been solved.

The results are not uniformly positive. On the hard eight-crossing unknot the
default full pipeline is about $1.38$ times slower even though raw backend
computation improves modestly. The extra invariant stages do not reject an
unknot and therefore add overhead. The small $T(3,5)$ and grid cases are also
slower because the original polynomial stage was already cheap. In absolute
terms these changes are small, but they are not omitted from the tables.

Three- and four-Conway sums finish quickly in the new recognizer, whereas all
baseline trials exhaust the selected resource budget. No completed-baseline
speedup is reported for these censored cases. On the selected 36-crossing
random five-braid, \emph{both raw backends exhaust the budget}. This remains an
unresolved performance case for the present code. Earlier timing claims in
the supplied archive are not substituted for fresh measurements.

\subsection{Counters that should not be compared naively}
The optimized \code{compositions} counter counts an explicit composition once
where it is actually called. The old counter increments by two inside each
inner elimination pair even though the first half-composition is outside that
inner loop. Their numeric values have different meanings; we do not use their
ratio as an operation-count speedup. Object counts, final ranks, and scan
orders retain their mathematical meanings. Bounded caches can also change
warm-process behavior; all headline data above are cold-process data.

\section{Validation and the limits of that validation}
The final command
\begin{lstlisting}
python -S -m unittest discover -s tests -v
\end{lstlisting}
passes \textbf{42 tests in 9.444 seconds} in the recorded run. This count is
unittest test methods, not 42 individual diagrams: many methods contain
exhaustive or randomized families. The source and full log are included.

\begin{itemize}[itemsep=3pt]
\item The original 19 tests are retained with narrowly scoped changes to
expected method names and to filter disabling in backend-resource tests.
The untouched original test file is preserved in \code{reference/}.
\item On 168 one-component three-braid words of lengths two and four, and
additional deterministic random samples, new verdicts are compared with an
independent reduced resolution cube. A separate 80-diagram test compares raw
rank and rank-by-cube-degree with the original backend, with $d^2=0$ checks.
\item The bracket is compared to an independent global state sum over exact
\code{Fraction} values on those 168 words, 120 random diagrams, and the
crossing-free case. Transition gluing is also compared to the separate
cobordism transition routine. Mirrors, arbitrary scan orders, R3 braid
relations, and positive/negative Markov stabilizations are tested.
\item All 2168 pairs of basis morphisms for triples of the five noncrossing
matchings on six endpoints are composed by both kernels. Another 250 random
eight-endpoint combinations test agreement and associativity. Every unit of
endomorphism rings through three arcs is checked, including its old inverse.
\item Every starting crossing is checked for planner equivalence on 100
random diagrams. Descending results match the old implementation on the
exhaustive family and 150 random diagrams. Writhe and Alexander behavior are
checked under random two-slot row rotations and row permutations.
\item Sparse determinants are compared with permutation-expansion determinants
for sizes zero through five. Modular Alexander minors are checked against
full exact polynomials on the exhaustive family and 200 random diagrams.
Deliberate modular collisions, budget exhaustion, malformed orders, invalid
diagrams, cache mutation, witness tampering, and command-line exit codes are
covered.
\item The independent reduced cube computes rank 33 for both Conway and
Kinoshita--Terasaka. Connected-sum tests verify the 22-crossing fixture, rank
1089, and exact integer bracket powers for the first five Conway sums.
\end{itemize}

Random tests use fixed seeds and include diagrams that may represent the same
knot; no claim about an exhaustive knot census is made. The independent cube
and bracket state sum avoid the optimized composition and gluing code, which
reduces the risk of a shared implementation error. The PD parser and ordinary
Python arithmetic remain shared trusted infrastructure. This is not a Lean
formalization, a proof of correctness of the Python interpreter, or exhaustive
verification for all inputs. The mathematical arguments explain the intended
algorithm; the tests substantiate the implementation on the stated domains.

\clearpage
\section{Using the package and interpreting evidence}
\subsection{Commands}
Run from the extracted package root; installation is optional.
\begin{lstlisting}
python -m fastunknot recognize examples/conway.json
python -m fastunknot recognize examples/conway_sum_2.json --output result.json
python -m fastunknot verify examples/conway_sum_2.json result.json

# Compute homology itself, not just a knottedness obstruction.
python -m fastunknot khovanov examples/conway.json --check-d2

# Force the complete backend; no preprocessing or invariant shortcut.
python -m fastunknot recognize examples/hard_unknot_8.json \
  --no-reduction --no-descending --no-alexander --no-jones --check-d2

# An exact one-point obstruction, not a complete recognizer.
python -m fastunknot bracket examples/conway.json --integer

# Reproduce and regenerate the recorded tables.
python benchmarks/run.py --repeats 3 --seconds 3
python benchmarks/make_tables.py
cd docs
pdflatex -interaction=nonstopmode -halt-on-error unknot_acceleration.tex
pdflatex -interaction=nonstopmode -halt-on-error unknot_acceleration.tex
\end{lstlisting}

The Python API retains \code{Diagram}, \code{recognize}, and
\code{khovanov_rank}. A supplied raw scan order must be a permutation of all
crossing indices; duplicate, missing, Boolean, and out-of-range indices are
rejected. Directly constructed \code{Diagram} objects are revalidated by the
recognizer. The default result JSON explicitly sets
\code{quasipolynomial_guarantee} to false.

\subsection{Limits and ablation switches}
\code{--no-alexander} disables both the modular determinant and full Alexander
polynomial; it does not disable the independent bracket tests.
\code{--no-modular} disables just the modular determinant, not the bracket.
\code{--no-jones} disables both bracket evaluations.
\code{--no-integer-jones} disables only the exact integer follow-up.
Setting the transition budget to zero forces the bracket stage to pass through
without a decision. The low-level Python bracket function accepts a custom
order and \code{modulus=None} for exact integer mode; omitting its caps gives
an uncapped filter computation.

\code{--max-objects} caps the Khovanov complex, not the invariant filters. Thus
a knottedness obstruction may succeed even with \code{--max-objects 1}.
\code{--seconds} is cooperative. Checks occur around stages and within
elimination loops, but a large single operation, validation, planning, or
Reidemeister simplification can overshoot it. A strict wall-clock or memory
policy requires an operating-system process limit, as in the benchmark runner.
Limits should never be interpreted as a complexity guarantee.

Recognition exit status is zero for either exact verdict, two for invalid
input, and three for \code{UNKNOWN}. Verification returns zero for accepted
evidence and four for rejected evidence. Successful \code{bracket} execution
returns zero even when \code{detected} is false; inspect its JSON fields rather
than interpreting this utility's exit status as a knot verdict.

\subsection{What the verifier certifies}
A modular or integer \code{KNOTTED} result contains its evaluation parameters,
computed value, required unknot value, scan order where relevant, and the
R1/R2 reduction trace. \code{verify_modular_result} first replays legal moves
on the supplied original input, then recomputes the obstruction on the reduced
diagram. Its historical function name includes ``modular'', but it also
supports exact integer bracket evidence.

For the Conway example at $A=2$ and $p=2147483647$, the recorded writhe is $-1$,
the bracket residue is $142097288$, and the required unknot residue is
$1140850688$. They differ. The corresponding result and verification output
are in \code{results/witnesses/}. This is evidence about the original input,
not merely an unrelated simplified diagram.

The verifier does not accept arbitrary Khovanov-rank JSON as a proof, and it is
not a small-certificate or polynomial-time verifier claim: bracket verification
recomputes the scan and can be expensive. The record is a reproducible exact
obstruction, not a cryptographic attestation or an adversarially hardened
network protocol.

\section{Conclusions and remaining engineering opportunities}
The supplied implementation meets the practical-performance alternative in
the request, with a large completed same-task speedup, a faster backend,
proved asymptotic improvements in ordering and descending detection, and a
family-specific separation between exact bracket rejection and explicit
Khovanov materialization. It also exposes a concrete mistake in treating
frontier matching count as a bound on complex size.

The largest outstanding performance issue is inputs on which all inexpensive
obstructions are inconclusive and the minimal complex has large
multiplicities. The present code still constructs that complex. Potential next
engineering steps include factorization across proven connected-sum cuts,
compressed representations of repeated summands, incremental simplification,
and different elimination pivot policies. None is silently credited to this
implementation or to the measurements. A global quasi-polynomial result would
need a proved structural replacement, such as a fully constructive accelerated
hierarchy with controlled depth and coordinate cost. The code, proofs,
measurements, and limitations in this delivery are kept separate so that such
future progress can be assessed against a reproducible baseline.

\clearpage
\begin{thebibliography}{9}
\raggedright
\bibitem{Input}
\emph{Knots.zip}, user-supplied archive, inspected 18 September 2026.
Baseline: \code{fast/fastunknot/}; original tests:
\code{fast/tests/test_fastunknot.py}. Literature materials:
\code{docs/quasipolynomial-talk.pdf} and
\code{docs/arXiv-2607.23350v1/algorithm-incompressible-250726.tex}.

\bibitem{LackenbyTalk}
M.~Lackenby, \emph{Unknot recognition in quasi-polynomial time}, talk slides,
February 2021, 109 physical PDF pages. Official author copy:
\url{https://people.maths.ox.ac.uk/lackenby/quasipolynomial-talk.pdf}.

\bibitem{Lackenby2026}
M.~Lackenby, \emph{Incompressible surfaces, hierarchies and unknot recognition},
arXiv:2607.23350v1, 25 July 2026. In particular Section 9, pp.~34--35,
and the refinements in Sections 10--14.
\url{https://arxiv.org/abs/2607.23350}.

\bibitem{BN}
D.~Bar--Natan, \emph{Fast Khovanov homology computations},
Journal of Knot Theory and Its Ramifications \textbf{16} (2007), 243--255.
\url{https://arxiv.org/abs/math/0606318}.

\bibitem{Kauffman}
L.~H.~Kauffman, \emph{State models and the Jones polynomial},
Topology \textbf{26} (1987), 395--407.
DOI: \href{https://doi.org/10.1016/0040-9383(87)90009-7}
{10.1016/0040-9383(87)90009-7}.

\bibitem{KM}
P.~B.~Kronheimer and T.~S.~Mrowka,
\emph{Khovanov homology is an unknot-detector},
Publications Math\'ematiques de l'IH\'ES \textbf{113} (2011), 97--208.
\url{https://arxiv.org/abs/1005.4346}.

\bibitem{Shumakovitch}
A.~Shumakovitch, \emph{Torsion of the Khovanov homology},
Fundamenta Mathematicae \textbf{225} (2014), 343--364;
arXiv revision dated 19 June 2018.
\url{https://arxiv.org/abs/math/0405474}.

\bibitem{KhovanovPatterns}
M.~Khovanov, \emph{Patterns in knot cohomology I},
Experimental Mathematics \textbf{12} (2003), 365--374.
\url{https://arxiv.org/abs/math/0201306}.
\end{thebibliography}
\end{document}
